\documentclass[12pt]{unisenaiclass}
\usepackage{unisenaistyle}
\usepackage[utf8]{inputenc}
\usepackage[brazil]{babel}
\usepackage{graphicx}
\usepackage{float}
\usepackage{amsmath}
\usepackage{hyperref}
\usepackage{caption}

\usepackage{csquotes} % Recomendado com biblatex
\usepackage[backend=biber, style=abnt]{biblatex} % ou 'style=apa', 'numeric', etc.
\addbibresource{referencias.bib} % arquivo .bib com as referências

\title{TÍTULO DO ARTIGO}
\subtitle{ARTICLE TITLE}

\author{
    \textbf{Gizela Alves Pieri Santos } \thanks{}\\
    \textbf{} \thanks{}
}

\begin{document}

\maketitle

\begin{resumo}
Esse é a seção do resumo
\end{resumo}

\begin{justify}
    \textbf{Palavras-chave:}  IoT; Backend; API; Sensores; Automação.
\end{justify}

\begin{abstract1}
This paper presents the development of an IoT hub integrated with a Node.js backend and a Next.js web interface. The system enables monitoring of temperature, humidity, pressure, and presence sensors, as well as remote control of an actuator (relay). The methodology involved the creation of a RESTful API to ensure communication between frontend and backend, providing scalability and interoperability. Results demonstrate the feasibility of the solution for residential, industrial, and agricultural automation applications.
\end{abstract1}

\begin{justify}
    \textbf{Keywords:}IoT; Backend; API; Sensors; Automation.  
\end{justify}



\section{Introdução}
A Internet das Coisas (IoT) tem se consolidado como uma tecnologia essencial para a automação e monitoramento de dispositivos. Este trabalho busca apresentar um sistema capaz de integrar sensores e atuadores a uma plataforma web, permitindo controle remoto e análise em tempo real. O objetivo é demonstrar a viabilidade técnica da solução e suas possíveis aplicações em diferentes contextos.
\section{Revisão de Literatura}
Segundo Aguiar (2023), a Internet das Coisas (IoT) pode ser descrita como uma rede de objetos físicos conectados à internet, capazes de interagir por meio de sensores e softwares. O autor destaca que a utilização de APIs RESTful é fundamental para conectar o mundo físico ao digital, permitindo que dispositivos troquem dados de forma eficiente e escalável.

De acordo com "Rodrigues (2024), a expansão de funcionalidades em sistemas IoT depende diretamente da integração com APIs HTTP, que possibilitam maior flexibilidade e controle remoto dos dispositivos. A autora enfatiza que o uso de arquiteturas REST garante interoperabilidade entre diferentes plataformas e tecnologias, tornando os sistemas mais adaptáveis às necessidades dos usuários".

A literatura recente sobre IoT aponta que a integração com APIs RESTful é essencial para garantir interoperabilidade e escalabilidade em ambientes complexos, permitindo que sensores e atuadores sejam facilmente conectados a aplicações web e móveis (Aguiar, 2023; Rodrigues, 2024).

\begin{itemize}
    \item \textbf{Exemplo de citação direta longa:}
\end{itemize}

Segundo \textcite{Segundo Rodrigo (2024)}, A literatura recente reforça que a integração entre IoT e APIs RESTful é essencial para garantir interoperabilidade, escalabilidade e confiabilidade em ambientes complexos.

\begin{quote}
Citação direta curta:  
Segundo Rodrigo (2024), “a interoperabilidade em IoT só é possível por meio de APIs RESTful bem definidas”.

Citação direta longa:  
Rodrigo (2024) afirma:

“A utilização de APIs RESTful em sistemas IoT não apenas garante a interoperabilidade entre dispositivos, mas também possibilita a escalabilidade e a segurança necessárias para aplicações críticas. Sem esse padrão de comunicação, os sistemas tornam-se fragmentados e pouco confiáveis.”

Citação indireta:  
De acordo com Rodrigo (2024), a integração de APIs RESTful em sistemas IoT é fundamental para assegurar interoperabilidade e escalabilidade, permitindo que diferentes dispositivos troquem informações de maneira segura e eficiente.
   
\end{quote}

Por outro lado, a citação indireta refere-se à reformulação das ideias do autor com as próprias palavras do estudante, sem copiar literalmente o texto. Mesmo assim, é obrigatório indicar a fonte. Por exemplo: O sistema TeX foi desenvolvido por Knuth com o objetivo de produzir livros com alta qualidade tipográfica \cite{knuth1984tex}.

Ao final do documento, as referências serão listadas automaticamente.

\section{Metodologia}
O sistema foi desenvolvido em duas camadas:
Backend Node.js: responsável pela criação da API RESTful e gerenciamento dos sensores e atuadores.

Frontend Next.js: interface gráfica para monitoramento e controle, utilizando React e hooks para atualização periódica dos dados.





\section{Resultados e Discussões}
O sistema apresentou:

Monitoramento em tempo real dos sensores.

Controle remoto do relé.

Alternância do status de conexão do dispositivo
\subsection{Exemplo de Tabela}


\begin{table}[H]
    \centering
    \caption{ Monitoramento de temperatura, umidade, pressão e presença.}
    \label{tab:exemplo}
    \begin{tabular}{|c|c|c|}
        \hline
        Sensor & Valor Atual & 	Unidade\\ \hline
        Temperatura  & 25 & °C \\ \hline
        Umidade& 60 &  porc \\ \hline
         Pressão & 1013 &  hPa \\ \hline
         Presença & Detectada&  -\\ \hline
         
    
    \end{tabular}
\end{table}

\section{Conclusão}
O projeto demonstrou a viabilidade de integrar sensores e atuadores em uma plataforma IoT com backend e API. Futuramente, pretende-se expandir a solução para incluir análise preditiva e integração com serviços em nuvem
\section*{Agradecimentos}
Agradeço à instituição de ensino e ao orientador Gabriel que contribuiu para o desenvolvimento deste trabalho.


\printbibliography
AGUIAR, Claudio. API REST e Internet das Coisas: Conectando o mundo físico ao digital. DIO, 2023.

RODRIGUES, Thaís dos Santos. Expansão de funcionalidades em sistema IoT de controle de acesso através de API HTTP. Universidade Federal de Pernambuco, 2024.
\end{document}